\chapter{Conclusions and Future Work}
\epigraph{``I don't see that human intelligence is something that humans can never understand."''}{\vspace{10pt}John McCarthy, March 1989}
\label{chapter:conclusions-fw}

\newpage

This chapter presents the final conclusions obtained through the research effort and the future work that can be derived. 



\section{Conclusions}

This project focused on evaluating the impact of modifying the MLP head in a Visual Transformer-based Facial Expression Recognition model specifically using a challenging or \gls{in-the-wild} based Datased and using \gls{POSTER} as the baseline architecture without changes in the MLP classification head \cite{zheng_poster_2022}. This allows to have a different perspective regarding how the Visual Transformers can help achieve a better performance when using an enhanced version of the classified head for \gls{in-the-wild} Facial Expression Recognition (\gls{FER}). 

Thought the work developed in the thesis, a new version of the head is presented as an opportunity to explore how the final classification can help improve model predictability or also achieve better results on challenging classes without modifying the visual transformer's backbone. The new head was integrated into the code  using the same hyperparameters as the baseline and tested with three transformer's sizes (Base/Small/Large).

%
%\begin{itemize}
%	
%	\item  Conclusion 1: The experimental results confirm the research hypothesis that modifying the MLP head leads to improved FER performance compared to the baseline POSTER architecture.
%	
%	\item Conclusion 2.
%	
%	\item Conclusion 3.
%	
%	
%\end{itemize}

\textbf{Conclusion 1 - Head modifications can improve FER performance in a nuanced way.}


The experimental results partially confirm the research hypothesis. Across all six configurations, overall top‑1 accuracy remained stable $(\approx 0.923–0.928)$, indicating that the choice of head architecture does not drastically affect dataset-level accuracy. However, macro‑F1 and per-class metrics reveal that the enhanced head delivers consistent gains for specific emotions, especially in the Small and Base \gls{POSTER} variants. Thus, while global accuracy is largely preserved, the modified head improves the quality of predictions at a finer granularity, particularly for less represented or more ambiguous emotion classes (e.g. Fear, Sad and Disgust)~\cite{kollias_affect_2021, zhang_leave_2023, wang_survey_2024}.

\textbf{Conclusion 2 - The enhanced head benefits minority and difficult emotion classes.}

Detailed analysis of class-wise precision, recall, and F1‑Score shows that the enhanced head improves performance for challenging categories such as Fear, which are both underrepresented and harder to discriminate. These gains are achieved without harming, and in some cases slightly improving, well-defined majority classes like Happy~\cite{kollias_affect_2021, zhang_leave_2023, wang_survey_2024}. This suggests that the combination of non-linear transformation, normalization, and regularization in the head helps the model exploit subtle facial cues that a purely linear MLP head fails to capture, thereby increasing robustness to \gls{intra-class} variability and class imbalance.


\textbf{Conclusion 3 - Smaller \gls{POSTER} variants with enhanced heads offer an efficient accuracy–cost trade‑off.}


The comparison between Small, Base, and Large \gls{POSTER} models indicates that larger backbones do not yield substantial accuracy gains over more compact configurations when the classification head is properly designed. Training-time analysis shows that Small and Base models are significantly more efficient, while achieving almost identical accuracy and macro‑F1 as the Large model. Consequently, from a practical perspective, deploying a Small or Base \gls{POSTER} model with the enhanced head is preferable for many \gls{in-the-wild} \gls{FER} applications, as it provides competitive performance with lower computational cost and resource requirements.


\section{Future Work}

The current investigation allows to have a lot opportunities to outperform the \gls{POSTER} \cite{zheng_poster_2022} research head investigation 


\begin{itemize}
\item A usage of another dataset for validation is recommended to allow review head changes performance on the proposed head. This permits to have a more valid experiment and metrics for a proper benchmarking comparisson. The dataset Affecnet\footnote{\url{https://www.mohammadmahoor.com/pages/databases/affectnet/}} \footnote{\url{https://www.kaggle.com/datasets/fatihkgg/affectnet-yolo-format}} is used on the baseline architecture paper ~\cite{zheng_poster_2022}. 

Although RAF-DB was selected as the primary dataset for the thesis work due to its balanced annotations and widespread adoption in transformer-based \gls{FER} researches and state of the art papers, other commonly used datasets include FER2013 \footnote{\url{https://www.kaggle.com/datasets/msambare/fer2013}}~\cite{khaireddin_facial_2021} and AffectNet~\cite{mollahosseini_affectnet_2019}, which are often employed for large-scale or cross-dataset generalization studies.

\item  The enhanced head architecture explored in this work represents only one point in a broader design space. Future work could investigate deeper or multi‑branch heads, gated or adaptive activation functions (e.g., GLU‑based layers), residual MLP blocks, or attention‑augmented classification heads. Systematic ablation studies comparing these variants would help isolate which components (normalization, dropout, depth, activation type) most strongly contribute to performance gains on difficult emotion classes.

\item Whereas this thesis work relied on standard evaluation metrics including accuracy, macro‑averaged F1‑Score, precision, recall, confusion matrices, and learning curves and non parametric tests like McNemar and BootstrapAnalysis  additional statistical validation techniques can also explored like for example Cohen’s kappa for agreement analysis, and Chi‑Square tests on confusion matrices. The use of such methods provides richer information about the significance and reliability of observed differences. A valuable line of future work is to formalize and integrate these statistical tests into the evaluation pipeline encouraging more statistically rigorous benchmarking practices in \gls{FER}.
\end{itemize}. 



%\begin{itemize}
%	\item Future work 1.
%	
%	\item Future work 2.
%	
%	\item Future work 3.
%	
%\end{itemize}
